\section{Patched-Kernel Validation}
\label{sec:patched-kernel-validation}

Chapter~8 established that both Canonical's temporary module-loading rule and
the custom MORI controls are transitional measures. The stronger end state is
not a more elaborate denial policy, but a corrected kernel that permits the
affected interface to remain available without reproducing the page-cache
manipulation observed on the vulnerable runtime. Phase~05 therefore removes the
temporary controls from the intended decision path and returns to the original
proof-of-concept question:

\begin{quote}
\textbf{RQ6 -- Patch Validation.} Does reproduction fail after installation of
a corrected kernel, and do the observations from the final retest match the
expected behaviour of the upstream fix?
\end{quote}

The final retest used Ubuntu kernel \code{6.8.0-137-generic}, supplied by
package version \code{6.8.0-137.137}. Canonical identifies
\code{6.8.0-117.117} as the first fixed Noble generic-kernel version, placing
the tested package beyond the published correction boundary
\cite{ubuntu-copyfail-cve}. The ordinary \code{algif\_aead} loading path was
restored, the module was successfully loaded, and the tested benign AF\_ALG
AEAD operations remained available.

Against that state, the same Python sample hash that produced a UID~0 shell on
\code{6.8.0-116-generic} instead reached the normal \code{/usr/bin/su}
authentication boundary. An independently compiled C implementation produced
the same patched-side outcome in two recorded executions. The mounted view of
\code{/usr/bin/su} retained its known-good digest and the expected root-owned
SUID metadata throughout the captured validation.

These observations support a positive but bounded answer to RQ6. Reproduction
failed for the two tested implementations, and the combination of retained
interface availability with no captured privileged-file change is consistent
with the upstream return to out-of-place AEAD processing. The evidence does not
constitute a binary-level proof that a particular commit executed internally,
nor is the Phase~05 sequence a perfectly matched controlled trial. This chapter
therefore documents both the result and the provenance imperfections that limit
its causal strength.

\subsection{Validation Logic and Evidence Boundary}

Patch validation requires more than observing that a proof of concept prints
an error or fails to open a root shell. The same visible outcome could be
produced by a disabled module, a still-active compensating control, an altered
test sample, a different target state, or an unrelated execution failure. The
Phase~05 design consequently uses a chain of controls rather than treating the
proof-of-concept output as a self-sufficient result.

The required observations are summarised in
\cref{tab:patch-validation-criteria}.

\begin{table}[H]
    \centering
    \caption{Success criteria for the patched-kernel validation.}
    \label{tab:patch-validation-criteria}
    \small
    \begin{tabularx}{\textwidth}{@{}>{\bfseries\raggedright\arraybackslash}p{0.23\textwidth}>{\raggedright\arraybackslash}p{0.31\textwidth}Y@{}}
        \toprule
        Criterion & Required observation & Evidential meaning \\
        \midrule
        Corrected runtime identified &
        Running kernel and installed package are recorded at or beyond the
        vendor's fixed-version boundary. &
        Connects the retest to a vendor-remediated release rather than merely
        to a newer-looking version string. \\

        Alternative controls removed &
        No active MORI enforcement path is observed, and normal module
        resolution is restored. &
        Reduces the risk of attributing a custom or module-level denial to the
        kernel correction. \\

        Interface retained &
        \code{algif\_aead} can load and an unprivileged control reaches the
        tested AF\_ALG AEAD operations. &
        Distinguishes correction of the vulnerable behaviour from removal of
        the complete interface. \\

        Reproduction absent &
        The proof of concept does not produce the previously observed UID~0
        transition. &
        Demonstrates regression for the tested execution path. \\

        Target preserved &
        The mounted target retains its known-good digest, ownership, size, and
        SUID mode after the retests. &
        Excludes the successful changed-view state documented in Chapter~5 at
        the time of the post-test captures. \\

        Implementation diversity &
        A separately implemented and hashed C path is also exercised. &
        Reduces dependence on behaviour unique to the Python wrapper, without
        proving every exploit variant. \\
        \bottomrule
    \end{tabularx}
\end{table}

The frozen public package contains five text artifacts and thirteen
screenshots. Its \code{PUBLIC-SHA256SUMS.txt} manifest covers all eighteen
objects, and all entries re-verify against the preserved repository copy. The
README remains outside that manifest so that documentation can change without
being confused with captured experimental evidence. Conclusions in this
chapter are therefore grounded in the manifested objects rather than in the
README's descriptions alone.

The evidence classes are not interchangeable. Machine-readable artifacts
provide exact captured text and timestamps where present. Screenshots preserve
interactive executions that were not also recorded as raw transcripts. The
vendor record establishes the fixed package boundary. The upstream change
describes the expected implementation-level correction. No single class proves
the whole causal claim; the conclusion depends on their agreement.

\subsubsection{Known imperfections carried forward}

Four boundaries materially affect the evaluation:

\begin{enumerate}
    \item The Phase~05 baseline artifact carrying the \emph{pre-patch} label is
    misnamed. Its own contents identify \code{6.8.0-137-generic}; it is a
    patched pre-retest baseline, not a capture from before installation of the
    fix.

    \item The text artifact named
    \code{phase05-independent-c-poc-provenance.txt} contains permission errors
    where source and binary hashes were intended. The successful hash values
    exist in screenshot~10 and are attributed only to that screenshot.

    \item The machine-readable MORI-isolation record is timestamped after the
    timestamp visible in the C provenance screenshot. Screenshot~02 shows the
    required isolation state but contains no terminal timestamp. The filename
    sequence and surrounding documentation support the intended order, but
    timestamps alone do not establish strict pre-test causality.

    \item The Python comparison changes the executing account. The vulnerable
    run used \code{labuser} at UID~1001, while the patched screenshot records
    \code{researcher} at UID~1000. The sample hash is matched, but the complete
    execution context is not.
\end{enumerate}

These imperfections do not erase the observed patched-side result. They do
prevent the chapter from claiming a laboratory-perfect single-variable trial.

\subsection{Expected Behaviour of the Kernel Correction}

The upstream correction is titled \emph{crypto: algif\_aead -- Revert to
operating out-of-place}. It largely reverts the earlier in-place implementation
while retaining explicit copying of the associated data. The commit explains
that the source and destination originate from different mappings and changes
the AEAD request so that the transmitted scatter-gather list is used as the
cryptographic source while the receive scatter-gather list remains the
destination \cite{linux-copyfail-fix}.

For this laboratory, that change leads to three observable expectations:

\begin{enumerate}
    \item \textbf{The interface should remain present.} The correction changes
    buffer handling; it is not the temporary policy that prevents
    \code{algif\_aead} from loading.

    \item \textbf{Benign operations should remain reachable.} An unprivileged
    process should still be able to create, bind, and accept the tested AEAD
    socket when the ordinary module path is available.

    \item \textbf{The exploit's shared-page effect should disappear.} The
    tested attempt should no longer provide the changed privileged-file view
    that previously allowed the SUID execution path to produce UID~0.
\end{enumerate}

The third expectation is phrased at the observation level. Phase~05 did not
instrument the internal scatter-gather lists or trace the precise source and
destination pages used by the corrected request. Agreement between the
published patch mechanism and the external result is therefore consistency
evidence, not a direct runtime trace of the upstream code path.

Canonical's Noble record supplies the distribution boundary needed for that
comparison. Generic kernel \code{6.8.0-117.117} is listed as fixed, while the
laboratory ran the later \code{6.8.0-137.137} package
\cite{ubuntu-copyfail-cve}. Version ordering supports the presence of the
correction in the vendor build; the project did not independently extract and
compare the running kernel's machine code with upstream commit
\code{a664bf3d603dc3bdcf9ae47cc21e0daec706d7a5}.

\subsection{Patched Baseline and Control Isolation}

The first Phase~05 capture records the following environment:

\begin{lstlisting}
Operating System: Ubuntu 24.04.2 LTS
Kernel:           Linux 6.8.0-137-generic

linux-image-6.8.0-137-generic          6.8.0-137.137
linux-image-generic                    6.8.0-137.137

/usr/bin/su:
-rwsr-xr-x 1 root root 55680
c74311fe5636b7d7f9a56239fa8adeeab12ba86fe7d41b91afa85bf9bbdae78b
\end{lstlisting}

The older \code{6.8.0-116.116} packages remained installed, but
\code{uname} identified \code{6.8.0-137-generic} as the running release. This
distinction matters because installed packages do not establish which kernel
is executing.

The baseline also shows that both module-loading rules had been renamed with
disabled suffixes:

\begin{lstlisting}
disable-algif_aead.conf.disabled
mori-copyfail-guard.conf.disabled-v2-2-test
\end{lstlisting}

Screenshot~02 confirms the effective resolution rather than relying only on
those filenames. A dry run selected the normal \code{af\_alg.ko.zst} and
\code{algif\_aead.ko.zst} insertion paths. The same screenshot contains no MORI
process and no MORI BPF pin, and reports the integrity service as inactive. A
later machine-readable isolation record reports the service as both inactive
and disabled.

One distinction is necessary here. The earliest patched baseline still lists
the legacy \code{mori\_integrity\_watch.py} process and an enabled integrity
service. That component observes the protected target; it is not the MORI v2.7
BPF enforcement path. Nevertheless, the intended experiment removed even that
observer before the active retests so that its output could not be mistaken for
the kernel's decision. The later captures document the resulting inactive
state, subject to the timestamp limitation already stated.

\begin{table}[H]
    \centering
    \caption{State used to reduce alternative explanations in Phase~05.}
    \label{tab:phase05-isolation-state}
    \small
    \begin{tabularx}{\textwidth}{@{}>{\bfseries\raggedright\arraybackslash}p{0.25\textwidth}>{\raggedright\arraybackslash}p{0.31\textwidth}Y@{}}
        \toprule
        State question & Captured observation & Interpretation \\
        \midrule
        Which kernel ran? &
        \code{uname} and package capture identify
        \code{6.8.0-137-generic} / \code{6.8.0-137.137}. &
        Retest executed beyond Canonical's first fixed Noble version. \\

        Was vendor module denial active? &
        Rule stored under a disabled filename; dry run resolves to ordinary
        module insertion. &
        The temporary \code{/bin/false} policy was not the observed failure
        mechanism. \\

        Was MORI v2.7 active? &
        No MORI process or BPF pin appears in the isolation screenshot or the
        later text capture. &
        No captured custom enforcement path explains the patched result. \\

        Was \code{algif\_aead} available? &
        Normal resolution was restored and the module later appears loaded. &
        Result was not produced by wholesale interface removal. \\

        Is ordering airtight? &
        No; the untimestamped screenshot and later raw isolation capture do not
        prove the entire order from timestamps alone. &
        Causal attribution remains bounded rather than absolute. \\
        \bottomrule
    \end{tabularx}
\end{table}

\subsection{Benign AF\_ALG Regression Control}

Before evaluating proof-of-concept outcomes, Phase~05 tested whether an
ordinary unprivileged caller could still reach the affected interface. The
machine-readable control identifies \code{labuser} as UID~1001 and records:

\begin{lstlisting}
socket(AF_ALG):                  OK
bind(("aead", "gcm(aes)")):     OK
accept():                       OK

algif_aead             12288  0
af_alg                 32768  1 algif_aead
\end{lstlisting}

The target digest immediately following the control remained:

\begin{lstlisting}
c74311fe5636b7d7f9a56239fa8adeeab12ba86fe7d41b91afa85bf9bbdae78b
\end{lstlisting}

This control closes an important alternative explanation. A failed exploit on
a system where \code{algif\_aead} cannot load would reproduce the broad
temporary-mitigation result from Chapter~8, not demonstrate corrected in-kernel
behaviour. Here, the module was loaded and the three exercised operations were
available to the same unprivileged account class used by the original
reproduction.

The control is still limited. It does not perform a complete encryption and
decryption round trip, verify ciphertext or authentication-tag correctness,
measure performance, or exercise concurrent application workloads. Its
supported claim is \emph{tested interface availability}: socket creation, AEAD
binding to \code{gcm(aes)}, and \code{accept()} succeeded. The README's broader
phrase \emph{benign AEAD lifecycle} is interpreted only within those recorded
operations.

\subsection{Python Proof-of-Concept Differential Retest}

The Python path provides the strongest direct before-and-after comparison
because the sample digest is shared with the successful vulnerable-kernel run.

\subsubsection{Sample identity and execution-context difference}

Screenshot~07 records the patched-side source as a 731-byte Python file with
SHA-256:

\begin{lstlisting}
d401e7d1c00605749d6c617ace73ab20a762b72e41c2e1590331596e38219a61
\end{lstlisting}

Chapter~5 records the same digest for the Python proof of concept that produced
a UID~0 shell on \code{6.8.0-116-generic}. Matching the sample removes source
drift as an obvious explanation for the changed result.

The user context was not matched. The successful vulnerable run was executed by
\code{labuser}, UID~1001. The patched screenshot places the file under
\code{/home/researcher} and records the post-test process as
\code{researcher}, UID~1000. That account belongs to the local \code{sudo} and
\code{lxd} groups. The record does not show a \code{sudo} invocation and the
final identity remained UID~1000, but those supplementary memberships prevent
treating the account as equivalent to the deliberately restricted
\code{labuser} threat-model subject. This is a non-root result, not an
identity-perfect replay of Chapter~5.

\subsubsection{Observed result}

On the vulnerable kernel, the matched Python sample previously produced:

\begin{lstlisting}
uid=0(root) gid=1001(labuser) groups=1001(labuser),100(users)
root
\end{lstlisting}

On \code{6.8.0-137-generic}, the recorded execution instead reached:

\begin{lstlisting}
Password:
su: Authentication failure

uid=1000(researcher) gid=1000(researcher) ...
\end{lstlisting}

The authentication message is significant because the proof of concept
continued far enough to execute \code{/usr/bin/su}; the expected unauthorised
modification simply did not provide the prior bypass. It is not an explicit
kernel security-policy denial, and the chapter does not label it as one. The
patched observation is ordinary application behaviour after the attempted
manipulation failed to create the previously effective target view.

Screenshot~08 then records \code{/usr/bin/su} as a 55,680-byte root-owned file
with mode \code{4755} and the known-good digest. No changed mounted view was
captured after the Python attempt.

\subsection{Independent C Implementation Retest}

The second path was a separately compiled C implementation. It reduces the
chance that the patched-side observation is merely a Python runtime or wrapper
failure, but it does not provide a second matched vulnerable-versus-patched
pair: the preserved project record does not show this C binary succeeding on
the vulnerable kernel without MORI enforcement.

\subsubsection{Provenance evidence}

Screenshot~10 records the following source and binary identities:

\begin{lstlisting}
source  e1dec1348c0d43ab47d8f992d9a0336216fd37ff83e6fa1d9e39a80890fbd165
binary  3b2605dc5e820f40645f682385194469f9189cbb65ee99153b52aedf45be5706

gcc (Ubuntu 13.3.0-6ubuntu2~24.04.1) 13.3.0
\end{lstlisting}

The source and binary were owned by \code{labuser}. Screenshot~11 repeats the
binary hash immediately before the recorded executions and again between them.
This creates a useful visual provenance chain for the tested executable.

The raw artifact with the corresponding provenance filename does not contain
these values. Its \code{realpath}, \code{ls}, and \code{sha256sum} operations
all failed with \code{Permission denied}; only the compiler version succeeded.
The chapter therefore does not describe that text file as a successful hash
capture. The hashes above are supported by screenshots~10 and~11.

\subsubsection{Observed result}

The C implementation identified the selected \code{/usr/bin/su} offset,
reported ten attempted writes, and then executed the target. In both recorded
runs, the result was:

\begin{lstlisting}
[+] Executing /usr/bin/su
Password:
su: Authentication failure
\end{lstlisting}

The prompt returned to \code{labuser} after each execution. The program's own
line stating that it is ``Patching page cache'' describes its intended action;
it is not accepted as evidence that the privileged file view actually changed.
The independent evidence of effect is the target state and execution outcome,
both of which remained inconsistent with successful exploitation.

Screenshot~09 records the known-good target digest and metadata before the C
test. Screenshot~12 and the final machine-readable artifact later record the
same digest, root ownership, 55,680-byte size, and mode \code{4755}. A direct
invocation of \code{/usr/bin/su} also reached normal authentication and rejected
the unauthorised attempt.

\subsection{Differential Evaluation Against the Vulnerable State}

The most useful interpretation comes from comparing Phase~02 and Phase~05
rather than considering the patched screenshots in isolation. The differential
is summarised in \cref{tab:vulnerable-patched-differential}.

\begin{table}[H]
    \centering
    \caption{Observed vulnerable and patched states.}
    \label{tab:vulnerable-patched-differential}
    \scriptsize
    \begin{tabularx}{\textwidth}{@{}>{\bfseries\raggedright\arraybackslash}p{0.18\textwidth}>{\raggedright\arraybackslash}p{0.34\textwidth}Y@{}}
        \toprule
        Dimension & Vulnerable Phase~02 & Patched Phase~05 \\
        \midrule
        Running kernel &
        \code{6.8.0-116-generic}, package \code{6.8.0-116.116}. &
        \code{6.8.0-137-generic}, package \code{6.8.0-137.137}. \\

        Python sample &
        SHA-256 \code{d401e7d1...19a61}. &
        Same SHA-256 \code{d401e7d1...19a61}. \\

        Executing identity &
        \code{labuser}, UID~1001. &
        \code{researcher}, UID~1000; administrative supplementary groups; no
        \code{sudo} invocation captured. \\

        Module path &
        Normal \code{algif\_aead} insertion available and module loaded. &
        Vendor and MORI rules disabled; normal insertion restored; module
        loaded. \\

        Python outcome &
        UID~0 shell and username \code{root}. &
        Normal \code{su} authentication failure; identity remained UID~1000. \\

        Mounted target view &
        Digest changed to \code{44900c63...79be9} after exploitation. &
        Known-good digest \code{c74311fe...dae78b} retained in post-test
        captures. \\

        Backing-filesystem comparison &
        Direct ext4 read retained the known-good digest while the mounted view
        differed. &
        No equivalent direct backing-filesystem read was captured in Phase~05. \\

        Independent C path &
        No preserved unmitigated success trial for this binary. &
        Two executions reached normal authentication; final target state
        remained intact. \\
        \bottomrule
    \end{tabularx}
\end{table}

The paired Python result is the central regression evidence: the sample hash,
target path, operating-system family, module-dependent interface, and broad
execution objective remain the same, while the running kernel crosses the
vendor fix boundary and the privilege transition disappears. The different
account and incomplete timestamp chain mean the kernel version is not the only
changed recorded variable. The C result and benign control add converging
evidence, but they do not turn the sequence into a formally isolated causal
experiment.

The stable final target also needs careful interpretation. It demonstrates
that the altered mounted view found during the vulnerable experiment was not
present when the post-test hashes were captured. It does not exclude an
extremely short-lived intermediate change between measurements, because
Phase~05 did not continuously hash the file or trace every page-cache state.
The immediate normal \code{su} behaviour makes such an unobserved transient
insufficient for the tested privilege transition, but not logically impossible.

\subsection{Alignment with the Expected Upstream Behaviour}

RQ6 asks not only whether exploitation failed, but whether the result resembles
the expected consequence of the upstream change. The comparison is shown in
\cref{tab:upstream-observation-alignment}.

\begin{table}[H]
    \centering
    \caption{Alignment between patch expectations and laboratory observations.}
    \label{tab:upstream-observation-alignment}
    \scriptsize
    \begin{tabularx}{\textwidth}{@{}>{\bfseries\raggedright\arraybackslash}p{0.25\textwidth}>{\raggedright\arraybackslash}p{0.31\textwidth}Y@{}}
        \toprule
        Patch-level expectation & Laboratory observation & Strength of match \\
        \midrule
        AEAD source and destination handled out-of-place &
        The exploit no longer produced the effective changed target view used
        for UID~0. &
        Mechanistically consistent, but internal buffers were not traced. \\

        Associated data still copied and interface retained &
        \code{algif\_aead} loaded; socket, bind, and accept controls succeeded. &
        Direct support for tested interface availability, not complete crypto
        correctness. \\

        Kernel fix replaces need for module denial &
        Ordinary module resolution was restored while both PoC paths failed to
        produce elevated access. &
        Strong within the tested configuration, subject to MORI chronology
        caveat. \\

        Privileged target should remain unmodified by the exploit path &
        Mounted-view hash, ownership, size, and SUID mode matched the baseline
        after both retests. &
        Direct post-test evidence; no continuous or backing-device capture. \\
        \bottomrule
    \end{tabularx}
\end{table}

The observed pattern is more informative than a generic program failure. If
the temporary mitigation had remained active, AF\_ALG AEAD would have been
unavailable. If MORI v2.7 had made the decision, the expected result would have
been an explicit policy denial and associated MORI telemetry. Instead, the
tested interface remained reachable and the attempted exploit progressed to
ordinary \code{su} authentication without producing the prior target effect.
That pattern matches the external consequence expected from repairing the
in-place AEAD handling.

The wording remains \emph{matches the expected behaviour}, not \emph{proves the
upstream commit}. The installed package version, Canonical's fixed-status
record, and the matching external behaviour jointly support attribution to the
vendor kernel. Exact commit provenance inside the Ubuntu build and direct
runtime observation of the corrected scatter-gather path were not independently
captured.

\subsection{Limitations and Alternative Explanations}

The final claim is constrained by the following limitations:

\begin{itemize}
    \item \textbf{No single-variable snapshot pair.} The patched test followed
    several earlier experimental phases and was not produced by changing only
    the kernel inside one continuously recorded session.

    \item \textbf{Different Python user.} The same Python sample hash was used,
    but the vulnerable and patched runs used UID~1001 and UID~1000
    respectively.

    \item \textbf{Isolation chronology.} The state showing no MORI process or
    BPF pins is preserved, but its strict ordering relative to every test is not
    established by embedded timestamps alone.

    \item \textbf{Interactive-only execution evidence.} The successful Python
    and C patched retests are recorded in screenshots rather than complete raw
    terminal transcripts. No transcript was reconstructed retrospectively.

    \item \textbf{Failed machine-readable C provenance.} The text artifact
    contains permission errors; successful source and binary hashes are visual
    evidence from screenshot~10.

    \item \textbf{Limited benign control.} Socket creation, bind, and accept
    were tested. Encryption, decryption, authentication-tag verification,
    concurrency, and application-level compatibility were not.

    \item \textbf{No patched backing-device comparison.} Phase~05 did not
    repeat the direct ext4 read used in Chapter~5, so the differential relies on
    the mounted target state and execution outcome.

    \item \textbf{No universal exploit coverage.} Two implementations were
    exercised. The result does not prove that every variant, target, kernel
    flavour, or container-escape scenario is fixed in every deployment.

    \item \textbf{Version-based commit attribution.} The project relies on
    Canonical's version status and observed behaviour rather than an independent
    binary audit of the Ubuntu kernel package.
\end{itemize}

Several of these could be improved by a future evidence supplement: rename or
supersede the mislabelled baseline without deleting the original, repeat both
PoCs as the same dedicated unprivileged identity, capture timestamped isolation
immediately before each execution, record raw transcripts, repeat the backing
filesystem comparison, and perform a real AEAD encrypt/decrypt regression. Such
work would strengthen provenance and compatibility claims; it is not required
to rewrite the outcome already observed.

\subsection{Answer to RQ6}

RQ6 asks whether reproduction fails on a corrected kernel and whether the
observations match the upstream fix. For the tested environment and the two
preserved implementations, the answer is yes.

The same Python proof-of-concept hash that produced a UID~0 shell on Ubuntu
kernel \code{6.8.0-116-generic} did not reproduce that transition on
\code{6.8.0-137-generic}. It reached ordinary \code{/usr/bin/su}
authentication and the executing identity remained UID~1000. The independent C
binary likewise reached normal authentication in two recorded executions and
returned to the \code{labuser} prompt. After the retests, the mounted target
retained its known-good digest, root ownership, SUID mode, and normal
authentication behaviour.

The failure cannot be explained by simple removal of the tested interface:
\code{algif\_aead} loaded normally and the unprivileged socket, bind, and accept
control succeeded. No active MORI enforcement path appears in the isolation
state. The combination of retained interface availability and absent
privileged-file effect is consistent with the upstream correction's return to
out-of-place AEAD processing.

The supported conclusion is therefore:

\begin{lstlisting}
On Ubuntu 24.04.2 LTS with kernel 6.8.0-137.137, the tested
Python and independent C Copy Fail paths did not reproduce the
previous unauthorised privilege transition.

The tested AF_ALG AEAD interface operations remained available,
and /usr/bin/su retained its known-good captured state.

This outcome matches the expected external behaviour of the
upstream out-of-place correction, within the stated provenance
and experimental limitations.
\end{lstlisting}

This result closes the experimental arc: the vulnerable kernel reproduced the
effect, detection observed its consequence, temporary controls interrupted the
known path with different compatibility costs, and the corrected kernel removed
the reproduced behaviour without requiring continued module denial. The fixed
vendor kernel is therefore the appropriate final remediation for the tested
system.

\subsection{Transition to Final Synthesis}

The six research questions have now been answered individually. Chapter~10
draws them together rather than introducing another exploit experiment. It
separates what the laboratory demonstrated from what remains inferred,
evaluates the strengths and weaknesses of the evidence workflow, extracts the
defensive lessons from MORI's iterations, and identifies the work required
before any custom control could be considered outside a research setting.

The final synthesis must preserve the asymmetry visible throughout the report:
reproduction is specific, detection is behavioural but not uniquely
attributive, selective prevention is possible only within a bounded model, and
vendor remediation carries the strongest operational conclusion. Treating
those statements as equally broad would erase the most important lesson of the
project---assurance depends not only on whether a test passed, but on exactly
what the evidence was capable of testing.
